\documentclass[11pt,a4paper]{report}

% =========================
% Packages (minimum required)
% =========================
\usepackage[a4paper,margin=2.5cm]{geometry}
\usepackage{microtype}
\usepackage[T1]{fontenc}
\usepackage[utf8]{inputenc}
\usepackage[french]{babel}

\usepackage{amsmath,amssymb,amsthm,mathtools}
\usepackage{bm}
\usepackage{bbm}
\usepackage{graphicx}
\usepackage{booktabs}
\usepackage{longtable}
\usepackage{enumitem}
\usepackage{listings}
\usepackage{xcolor}
\usepackage{tikz}
\usepackage{pgfplots}
\pgfplotsset{compat=1.18}

\usepackage[hidelinks]{hyperref}
\usepackage[nameinlink,capitalise,noabbrev]{cleveref}

% =========================
% Colors / Listings
% =========================
\definecolor{codebg}{RGB}{248,248,248}
\definecolor{codeframe}{RGB}{210,210,210}
\definecolor{codetxt}{RGB}{20,20,20}
\definecolor{kw}{RGB}{0,92,170}
\definecolor{cm}{RGB}{90,110,90}
\definecolor{st}{RGB}{140,60,60}

\lstdefinestyle{py}{
language=Python,
basicstyle=\ttfamily\small\color{codetxt},
backgroundcolor=\color{codebg},
frame=single,
rulecolor=\color{codeframe},
numbers=left,
numberstyle=\ttfamily\tiny\color{black!60},
xleftmargin=2.2em,
framexleftmargin=1.8em,
showstringspaces=false,
tabsize=4,
breaklines=true,
keywordstyle=\color{kw}\bfseries,
commentstyle=\color{cm}\itshape,
stringstyle=\color{st},
upquote=true
}
\lstset{style=py}

% =========================
% Theorem-like environments
% =========================
\theoremstyle{definition}
\newtheorem{definition}{Définition}[chapter]
\newtheorem{example}{Exemple}[chapter]
\theoremstyle{plain}
\newtheorem{theorem}{Théorème}[chapter]
\newtheorem{proposition}{Proposition}[chapter]
\newtheorem{lemma}{Lemme}[chapter]
\theoremstyle{remark}
\newtheorem{remark}{Remarque}[chapter]

% =========================
% Macros (notations)
% =========================
\newcommand{\dd}{\mathrm{d}}
\newcommand{\E}{\mathbb{E}}
\newcommand{\Var}{\mathrm{Var}}
\newcommand{\Cov}{\mathrm{Cov}}
\newcommand{\Prob}{\mathbb{P}}
\newcommand{\1}{\mathbbm{1}}
\newcommand{\Normal}{\mathcal{N}}
\newcommand{\Pmeas}{\mathbb{P}}
\newcommand{\Qmeas}{\mathbb{Q}}
\newcommand{\R}{\mathbb{R}}
\newcommand{\F}{\mathcal{F}}
\newcommand{\argmin}{\operatorname*{arg,min}}
\newcommand{\argmax}{\operatorname*{arg,max}}

\newcommand{\filepath}[1]{\texttt{\detokenize{#1}}}
\newcommand{\codeid}[1]{\texttt{\detokenize{#1}}}
\newcommand{\NonImplem}{\textbf{non implémenté dans le code actuel}}

% =========================
% Boxes
% =========================
\newenvironment{jurybox}{
\par\medskip
\noindent\fcolorbox{black!30}{black!3}{\begin{minipage}{0.97\textwidth}
\textbf{Ce que le jury doit retenir}\par\smallskip
}{
\end{minipage}}
\par\medskip
}

\newcommand{\qa}[4]{
\begin{itemize}[leftmargin=*]
\item \textbf{À quoi ça sert ?} #1
\item \textbf{Hypothèses.} #2
\item \textbf{Limites.} #3
\item \textbf{Usage dans l'app.} #4
\end{itemize}
}

% =========================
% Document
% =========================
\begin{document}

% -------------------------
% Title page
% -------------------------
\begin{titlepage}
\centering
\vspace*{2cm}
{\Large\bfseries PaperTradingApp\par}
\vspace{0.4cm}
{\large Dossier théorique & cartographie d'implémentation (from ground up)\par}
\vspace{1.2cm}
{\large\itshape Public : jury financier (Master)\par}
\vspace{1.0cm}

\noindent\rule{0.9\textwidth}{0.4pt}\par
\vspace{0.6cm}

\begin{minipage}{0.9\textwidth}
\small
\textbf{Contraintes de ce dossier (rappel)}\par
\begin{itemize}[leftmargin=*]
\item Ce document est \textbf{strictement aligné} avec le code présent dans le dépôt \filepath{d:\PythonDProjects\PaperTradingApp}.
\item Toute affirmation sur \emph{l'implémentation} cite des \textbf{chemins de fichiers} et, quand pertinent, des \textbf{fonctions/classes} existantes.
\item Si une brique est mentionnée dans le plan mais absente du code, elle est marquée : \NonImplem.
\end{itemize}
\end{minipage}

\vfill
{\small Généré le \today\par}
\end{titlepage}

% -------------------------
% Abstract / toc / figures
% -------------------------
\chapter*{Résumé}
Cette application est une interface de \emph{paper trading} orientée «modèles» : elle consomme des données (marché, options, taux), construit des objets quantitatifs (surface de volatilité implicite, calibrations Heston/SABR/jumps/rough, courbe de taux, allocations), puis propose des décisions (pricing, couvertures, rebalancement) et peut exécuter des ordres via Alpaca (papier).

Le dossier est volontairement \textbf{pédagogique} : chaque grande fonctionnalité est reconstruite \textbf{à partir des hypothèses} (finance & probabilités), \textbf{dérivée} étape par étape, puis \textbf{reliée} à sa traduction dans le code. En particulier :
\begin{itemize}[leftmargin=*]
\item Pricing d'options : Black--Scholes (formule, grecs, volatilité implicite), solveur PDE (Crank--Nicolson), arbres CRR, Monte Carlo et Longstaff--Schwartz.
\item Volatilité : surface d'IV, interpolation/masques, calibration (moindres carrés) et modèles implémentés : SABR (Hagan), Merton (jumps), Bates, Heston (CF & FFT), rHeston (approximation markovienne), rBergomi (MC + surrogate), Volterra (MC proxy).
\item Couverture RL : un \emph{DQN} (numpy) entraîné sur environnement synthétique, puis utilisé comme aide à la couverture sur un proxy de delta.
\item Taux : courbe de taux continue, discount factors, forwards, calibration Nelson--Siegel (grid search sur $\tau$ + moindres carrés sur les $\beta$).
\item Allocation & risque : Markowitz (min-var / max Sharpe), risk parity (ERC), eigen-portfolio, métriques «risk management» (exposition, VaR-lite, alertes) et passage poids $\to$ ordres.
\end{itemize}

\tableofcontents
\listoffigures

% =========================================================
% CHAPITRE 0 — Architecture & Notations
% =========================================================
\setcounter{chapter}{-1}
\chapter{Architecture & Notations}

\section{A. Problème financier (objectif, inputs, outputs)}
L'objectif global est un pipeline \emph{data} $\to$ \emph{modèles} $\to$ \emph{décision} $\to$ \emph{exécution} :
\begin{itemize}[leftmargin=*]
\item \textbf{Inputs} : prix spot, historiques (OHLC), chaînes d'options (IV par strike/maturité), nœuds de courbe de taux, positions et compte Alpaca.
\item \textbf{Outputs} : prix d'options (et payoffs), grecs, surfaces d'IV (marché vs modèle), paramètres calibrés, suggestions de hedge (RL), plans de rebalancement, ordres soumis (paper).
\end{itemize}

\section{B. Théorie (définitions, hypothèses)}
\subsection*{Définition formelle (architecture logicielle)}
Le dépôt impose une séparation MVC stricte :
\begin{itemize}[leftmargin=*]
\item \textbf{Model} : logique métier pure (pricing, calibration, hedging, courbes, allocation). Répertoire : \filepath{app/model/}.
\item \textbf{Controller} : adaptateurs fins UI $\to$ Model. Répertoire : \filepath{app/controller/}.
\item \textbf{View} : UI Streamlit, rendu uniquement. Répertoire : \filepath{app/vue/}.
\item \textbf{Utils} : helpers sans logique métier. Répertoire : \filepath{app/utils/}.
\end{itemize}
Ces règles sont explicitement documentées dans \filepath{README.md} et \filepath{docs/architecture_mvc.md}.

\subsection*{Intuition (pourquoi c'est crucial en finance)}
Une application financière «modèles» doit pouvoir :
\begin{itemize}[leftmargin=*]
\item tester les moteurs de pricing (unit tests) sans UI ;
\item changer une source de data (Yahoo $\leftrightarrow$ Alpaca $\leftrightarrow$ Stooq) sans réécrire les modèles ;
\item tracer (audit) ce qui est calculé vs ce qui est affiché (séparation décision / présentation).
\end{itemize}

\qa
{Garantir qu'un résultat numérique (prix, calibration) est traçable, testable, et reproductible.}
{Discipline de dépendances : la Vue ne doit pas appeler le Model directement.}
{MVC ne garantit pas la qualité des modèles, seulement l'industrialisation et l'auditabilité.}
{Le point d'entrée Streamlit délègue à \filepath{app/vue/main_app.py} et les onglets appellent des contrôleurs \filepath{app/controller/.py} qui routent vers \filepath{app/model/.py}.}

\section{C. Dérivations (ultra détaillées)}
Ce chapitre est «ground up» au sens logiciel : on dérive le \emph{graphe de dépendances} (flux) observé dans le dépôt.

\subsection*{Cartographie repo (exigence de vérification)}
\paragraph{(1) Dossiers clés (racine).}
\begin{itemize}[leftmargin=*]
\item \filepath{app/} : code applicatif (MVC).
\item \filepath{cache/} : caches CSV (ex. \filepath{cache/YahooOptionChains/}, \filepath{cache/OHLC/}).
\item \filepath{data/} : snapshots JSON et données utilisateurs (ex. \filepath{data/yield_curves/}).
\item \filepath{docs/} : documentation d'architecture.
\item \filepath{scripts/} et \filepath{tests/} : intégrité MVC, tests pricing, etc.
\end{itemize}

\paragraph{(2) Points d'entrée (UI/CLI).}
\begin{itemize}[leftmargin=*]
\item UI Streamlit (Cloud) : \filepath{streamlit_app.py} appelle \codeid{app.vue.main_app.main}.
\item UI interne : \filepath{app/vue/main_app.py} construit la navigation et appelle les renderers des onglets.
\item CLI data (historique Stooq) : \filepath{app/model/market_data/scripts/fetch_history_cli.py} (\codeid{main}).
\end{itemize}
Remarque importante : \filepath{app/model/options/core/iv.py} tente d'appeler un script \filepath{app/scripts/fetch_history_cli.py} avec des arguments \codeid{--period/--interval} ; ce chemin/CLI ne correspond pas au script réellement présent (\filepath{app/model/market_data/scripts/fetch_history_cli.py} utilise \codeid{--start/--end/--freq}). On doit donc traiter cette voie comme \NonImplem\ au sens «non opérationnelle telle quelle».

\paragraph{(3) Modules métiers (Model).}
Exemples (non exhaustif) :
\begin{itemize}[leftmargin=*]
\item Options (pricing/grecs/PDE/MC) : \filepath{app/model/options/} (ex. \codeid{compute_price} dans \filepath{app/model/options/service.py}).
\item Volatility models & calibration : \filepath{app/model/volatility_models/} et \filepath{app/model/calibration/}.
\item Hedging RL (DQN) : \filepath{app/model/hedger_v2/}.
\item Yield curve : \filepath{app/model/yieldcurve/}.
\item Allocation & risk : \filepath{app/model/portfolio_allocation/}, \filepath{app/model/risk_management/}.
\item Market data : \filepath{app/model/market_data/market_data.py}.
\item Trading/exécution Alpaca : \filepath{app/model/alpaca_spot/}, \filepath{app/model/alpaca_orders/}.
\end{itemize}

\section{D. Numérique / algorithmes (stabilité, complexité, choix)}
Le dépôt mélange plusieurs familles d'algorithmes :
\begin{itemize}[leftmargin=*]
\item \textbf{Formules fermées} : Black--Scholes (ex. \filepath{app/model/options/engines/black_scholes.py}).
\item \textbf{PDE} : Crank--Nicolson pour BS en log-spot (classe \codeid{CrankNicolsonBS} dans \filepath{app/model/options/engines/pricing.py}).
\item \textbf{Arbres} : CRR binomial (ex. \codeid{price_american_crr} dans \filepath{app/model/options/engines/crr.py}).
\item \textbf{Monte Carlo} : simulateurs GBM et exotiques (ex. \codeid{simulate_gbm_paths} dans \filepath{app/model/options/core/trees.py}).
\item \textbf{FFT} : Carr--Madan pour prix de calls (ex. \codeid{carr_madan_fft_call_prices} dans \filepath{app/model/volatility_models/common/fft.py}).
\item \textbf{Optimisation} : moindres carrés multi-start (SciPy) et design LHS (ex. \codeid{latin_hypercube_samples} dans \filepath{app/model/calibration/optimizers.py}).
\item \textbf{RL} : DQN numpy (replay buffer + target network) dans \filepath{app/model/hedger_v2/dqn_hedger.py}.
\end{itemize}
Le choix «plusieurs algorithmes» est cohérent en finance : on cherche un compromis précision/temps et une robustesse (fallbacks).

\section{E. Implémentation dans le repo (fichiers, fonctions, pipeline, paramètres)}
\subsection*{Diagramme de flux (imposé)}
\begin{figure}[h]
\centering
\begin{tikzpicture}[
node distance=1.2cm and 1.2cm,
box/.style={draw, rounded corners, align=center, minimum width=3.6cm, minimum height=0.9cm},
arr/.style={-latex, thick}
]
\node[box] (md) {Market Data\\small (Stooq/Yahoo/Alpaca)};
\node[box, right=of md] (opt) {Options & Vol\\small (IV, pricing)};
\node[box, right=of opt] (cal) {Calibration\\small (SABR/Heston/...)};
\node[box, right=of cal] (hedge) {Hedging RL\\small (DQN)};
\node[box, below=of cal] (yc) {Yield Curve\\small (DF, forwards, NS)};
\node[box, left=of yc] (pa) {Portfolio Allocation\\small (Markowitz, RP, eigen)};
\node[box, left=of pa] (exec) {Trading/Execution\\small (Alpaca orders)};

\draw[arr] (md) -- (opt);
\draw[arr] (opt) -- (cal);
\draw[arr] (cal) -- (hedge);
\draw[arr] (yc) -- (pa);
\draw[arr] (pa) -- (exec);

\draw[arr, dashed] (yc) -- (opt);
\draw[arr, dashed] (cal) -- (opt);

\end{tikzpicture}
\caption{Flux logique cible : données $\to$ modèles $\to$ décision $\to$ exécution. Les flèches en pointillés correspondent à une intégration observable via \filepath{app/vue/components/options/router.py} (taux depuis Yield Curve) et \filepath{app/vue/tabs/tab_advanced_calibration.py} (envoi d'une surface calibrée vers Options).}
\label{fig:global-flow}
\end{figure}

\subsection*{Tableau de cartographie (module / rôle / fichiers / fonctions clés / I/O)}
\begin{longtable}{@{}p{3.0cm}p{3.6cm}p{4.8cm}p{4.2cm}@{}}
\toprule
\textbf{Module} & \textbf{Rôle} & \textbf{Fichiers (exemples)} & \textbf{Fonctions/classes clés (exemples)}\
\midrule
\endhead
Vue (UI) & Navigation + rendu Streamlit & \filepath{streamlit_app.py}, \filepath{app/vue/main_app.py}, \filepath{app/vue/tabs/tab_.py} & \codeid{main} (\filepath{app/vue/main_app.py}), \codeid{render_tab} (tabs)\
Contrôleurs & Glue View $\to$ Model & \filepath{app/controller/options_controller.py}, \filepath{app/controller/calibration_controller.py}, \filepath{app/controller/yieldcurve_controller.py}, \filepath{app/controller/hedger_v2_controller.py}, \filepath{app/controller/portfolio_and_risk_controller.py} & \codeid{run_advanced_surface_calibration}, \codeid{get_dqn_hedge_suggestion}, \codeid{compute_allocation}, \ldots\
Options (Model) & Pricing/grecs/PDE/MC/exotiques & \filepath{app/model/options/service.py}, \filepath{app/model/options/engines/pricing.py}, \filepath{app/model/options/engines/black_scholes.py}, \filepath{app/model/options/core/trees.py} & \codeid{compute_price}, \codeid{CrankNicolsonBS}, \codeid{black_scholes_price}, \codeid{simulate_gbm_paths}\
Calibration (Model) & Parsing surface, Heston V1, NN Heston & \filepath{app/model/calibration/market_surface.py}, \filepath{app/model/calibration/heston_calibrator.py}, \filepath{app/model/calibration/heston_nn.py} & \codeid{build_fixed_grid}, \codeid{calibrate_heston_least_squares}, \codeid{HestonSurfaceNet}\
Vol models (Model) & SABR/Heston/jumps/rough/Volterra & \filepath{app/model/volatility_models/} & \codeid{hagan_black_iv}, \codeid{carr_madan_fft_call_prices}, \codeid{rheston_log_return_cf_markovian}, \ldots\
Hedger RL (Model) & DQN numpy + envs & \filepath{app/model/hedger_v2/dqn_hedger.py}, \filepath{app/model/hedger_v2/env.py} & \codeid{DQNAgent}, \codeid{ReplayBuffer}, \codeid{SyntheticHedgingEnv}\
Yield curve (Model) & Courbe, DF, forwards, Nelson--Siegel & \filepath{app/model/yieldcurve/service.py}, \filepath{app/model/yieldcurve/engine.py}, \filepath{app/model/yieldcurve/rates_utils.py} & \codeid{get_active_curve}, \codeid{fit_nelson_siegel_grid_search}, \codeid{get_r}\
Allocation & risk (Model) & Markowitz/RP/eigen + VaR-lite & \filepath{app/model/portfolio_allocation/engine.py}, \filepath{app/model/risk_management/engine.py} & \codeid{markowitz_optimize}, \codeid{risk_parity_optimize}, \codeid{compute_var_lite}\
\bottomrule
\end{longtable}

\section{F. Exemple end-to-end (ce qui se passe quand l'utilisateur lance la fonctionnalité)}
\paragraph{Exemple 1 : calibrer une surface puis pricer dans Options.}
\begin{enumerate}[leftmargin=*]
\item L'utilisateur ouvre l'onglet «Calibration avancée» (\filepath{app/vue/tabs/tab_advanced_calibration.py}, \codeid{render_tab}).
\item Il charge une surface (Yahoo/CSV/Cache/Alpaca) puis lance la calibration via \codeid{CalibrationController.run_advanced_surface_calibration} (\filepath{app/controller/calibration_controller.py}).
\item Le contrôleur instancie un calibrateur (ex. \codeid{HestonFFTCalibrator} dans \filepath{app/model/volatility_models/heston/calibrator_fft.py}) et renvoie \codeid{iv_model} et \codeid{params}.
\item Un clic «Envoyer IV modèle vers Options» écrit \codeid{st.session_state["calib_model_surface_df"]} et force la source «Calibration» côté Options (\filepath{app/vue/tabs/tab_advanced_calibration.py}, bouton à proximité du bloc résultat).
\item L'onglet Options (\filepath{app/vue/components/options/router.py}, \codeid{render_options_router}) lit cette surface et l'affiche/consomme.
\end{enumerate}

\paragraph{Exemple 2 : couverture DQN.}
\begin{enumerate}[leftmargin=*]
\item L'utilisateur ouvre «Hedging Systems» (\filepath{app/vue/tabs/tab_hedging_systems.py}).
\item Il entraîne le modèle (bouton «Train / update DQN») : appel \codeid{ctrl.train_dqn_model} $\to$ \codeid{load_or_train_dqn_model} dans \filepath{app/model/hedger_v2/dqn_hedger.py}.
\item Il demande une suggestion : \codeid{ctrl.get_dqn_hedge_suggestion} $\to$ \codeid{suggest_hedge_action}.
\item Optionnellement, il exécute : \codeid{ctrl.execute_dqn_hedge} $\to$ \codeid{AlpacaHedgerClient.submit_market_order} (\filepath{app/model/hedger_v2/alpaca_client.py}).
\end{enumerate}

\section{G. Limites & extensions crédibles}
\begin{itemize}[leftmargin=*]
\item La présence de plusieurs moteurs (formule/PDE/tree/MC/FFT) exige une gouvernance : quels moteurs sont «référence» par produit ?
\item Certaines briques existent comme \emph{scaffold} : \codeid{app.model.calibration.logic.run_calibration} renvoie explicitement «Calibration non implémentée» (\filepath{app/model/calibration/logic.py}). L'onglet «Calibration avancée» utilise une voie séparée (contrôleur unifié) qui, elle, est opérationnelle.
\item Dividendes : \codeid{get_q} retourne 0.0 (placeholder) dans \filepath{app/model/yieldcurve/rates_utils.py} ; en production, il faut une source robuste (dividendes attendus/forwards).
\end{itemize}

\begin{jurybox}
L'application n'est pas «un modèle», mais une \textbf{boîte à outils} : elle assemble data, pricing, calibration, risque, exécution. Le point clé pour un jury est la \textbf{traçabilité} : chaque sortie observable dans l'UI correspond à une fonction repérable dans \filepath{app/model/} via \filepath{app/controller/}.
\end{jurybox}

% =========================================================
% PARTIE I — Stochastic calculus & pricing
% =========================================================
\part{Stochastic calculus & pricing (base indispensable)}

% ---------------------------------------------------------
\chapter{Brownien et intégrale stochastique}

\section{A. Problème financier (objectif, inputs, outputs)}
On veut modéliser l'incertitude du prix d'un sous-jacent (action) à court terme pour :
\begin{itemize}[leftmargin=*]
\item pricer des options (prix = espérance actualisée sous une mesure appropriée) ;
\item simuler des trajectoires (Monte Carlo) ;
\item construire des risques (grecs, PnL, mesures de risque).
\end{itemize}
\textbf{Input} : paramètres $(S_0,r,q,\sigma,T)$ et un générateur de bruit. \textbf{Output} : une trajectoire $S_t$ (ou une loi de $S_T$).

\section{B. Théorie (définitions, hypothèses)}
\subsection*{Définition formelle}
\begin{definition}[Mouvement brownien]
Sur un espace probabilisé filtré $(\Omega,\F,(\F_t){t\ge 0},\Pmeas)$, un processus $(W_t){t\ge 0}$ est un \emph{mouvement brownien standard} si :
\begin{enumerate}[leftmargin=*]
\item $W_0=0$ p.s. ;
\item $W_t-W_s \sim \Normal(0,t-s)$ pour $0\le s<t$ ;
\item les accroissements sont indépendants ;
\item $t\mapsto W_t$ est continu.
\end{enumerate}
\end{definition}

L'intégrale stochastique (Itô) est l'opérateur qui donne sens à
[
\int_0^T H_t,\dd W_t
]
pour des processus adaptés $H_t$ (typiquement carrés-intégrables).

\subsection*{Intuition financière}
\begin{itemize}[leftmargin=*]
\item Le brownien représente un «bruit» agrégé (micro-chocs) qui, à l'échelle du modèle, a des incréments gaussiens.
\item L'intégrale stochastique est la façon correcte d'accumuler un risque instantané : si l'exposition $H_t$ varie dans le temps, la volatilité cumulée n'est pas une somme classique mais une intégrale par rapport à $W$.
\end{itemize}

\qa
{C'est le socle de la dynamique GBM (Black--Scholes) et des simulations Monte Carlo utilisées partout (pricing, exotiques, calibration rough/Volterra).}
{Marché frictionless pour les modèles de base ; dynamique continue ; volatilité constante si GBM ; indépendance/gaussianité du bruit.}
{Le brownien est une idéalisation : sauts, volatilité stochastique, microstructure, queues épaisses.}
{Les simulateurs GBM dans \filepath{app/model/options/core/trees.py} et \filepath{app/model/options/logic.py} utilisent des incréments gaussiens (numpy) pour approximer $W$.}

\section{C. Dérivations (ultra détaillées)}
\subsection*{De $W_t$ à la dynamique de prix (GBM)}
Le modèle de Black--Scholes suppose une dynamique (sous la mesure risque-neutre $\Qmeas$) :
[
\frac{\dd S_t}{S_t}=(r-q),\dd t + \sigma,\dd W_t^{\Qmeas}.
]
On le réécrit :
[
\dd S_t = (r-q)S_t,\dd t + \sigma S_t,\dd W_t.
]
L'objet \emph{non classique} est $\dd W_t$ : ce n'est pas une dérivée ordinaire, mais un incrément brownien.

\subsection*{Solution exacte discrétisée}
On sait (via Itô, démontré au chapitre Itô) que la solution est :
[
S_t = S_0\exp\Big((r-q-\tfrac12\sigma^2)t+\sigma W_t\Big).
]
Pour une discrétisation $t_{n+1}=t_n+\Delta t$,
[
\ln S_{t_{n+1}}-\ln S_{t_n}=(r-q-\tfrac12\sigma^2)\Delta t+\sigma (W_{t_{n+1}}-W_{t_n}),
]
et comme $W_{t_{n+1}}-W_{t_n}\sim\Normal(0,\Delta t)$, on simule :
[
W_{t_{n+1}}-W_{t_n}=\sqrt{\Delta t},Z_n,\quad Z_n\sim\Normal(0,1),
]
d'où
[
S_{t_{n+1}} = S_{t_n}\exp\Big((r-q-\tfrac12\sigma^2)\Delta t+\sigma\sqrt{\Delta t},Z_n\Big).
]

\section{D. Numérique / algorithmes (stabilité, complexité, choix)}
\begin{itemize}[leftmargin=*]
\item La simulation via exponentielle est \textbf{stable} (garantit $S_t>0$).
\item Complexité : $O(N_{\text{paths}}\times N_{\text{steps}})$.
\item Le choix $Z\sim\Normal(0,1)$ est le pont concret entre théorie (brownien) et code (numpy RNG).
\end{itemize}

\section{E. Implémentation dans le repo (fichiers, fonctions, pipeline, paramètres)}
\subsection*{Traduction directe de la formule GBM}
Le simulateur le plus lisible est \codeid{simulate_gbm_paths} dans \filepath{app/model/options/core/trees.py}. Extrait (raccourci) :

\begin{lstlisting}[caption={Simulation GBM (extrait) : \filepath{app/model/options/core/trees.py}, fonction \codeid{simulate_gbm_paths}.},label={lst:gbm}]
def simulate_gbm_paths(S0, r, q, sigma, T, steps, n_paths, seed=None):
dt = float(T) / max(int(steps), 1)
drift = (float(r) - float(q) - 0.5 * float(sigma) ** 2) * dt
vol = float(sigma) * math.sqrt(dt)
rng = np.random.default_rng(seed)
Z = rng.standard_normal((int(n_paths), max(int(steps), 1)))
increments = drift + vol * Z
log_paths = np.cumsum(increments, axis=1)
S = float(S0) * np.exp(log_paths)
S = np.concatenate([np.full((int(n_paths), 1), float(S0)), S], axis=1)
return S
\end{lstlisting}

\paragraph{Lecture ligne à ligne (lien formule $\leftrightarrow$ code).}
\begin{itemize}[leftmargin=*]
\item \codeid{dt} implémente $\Delta t=T/N$.
\item \codeid{drift} implémente $(r-q-\frac12\sigma^2)\Delta t$.
\item \codeid{Z} est la matrice des $Z_n\sim\Normal(0,1)$.
\item \codeid{increments = drift + vol * Z} implémente $(r-q-\frac12\sigma^2)\Delta t+\sigma\sqrt{\Delta t},Z_n$.
\item \codeid{log_paths = cumsum(increments)} somme les incréments de $\ln S$.
\item \codeid{S = S0 * exp(log_paths)} revient à $S=\exp(\ln S)$.
\end{itemize}

\subsection*{Où l'app l'utilise}
\begin{itemize}[leftmargin=*]
\item Options exotiques MC (Asian, lookback, barrières, cliquet) : \filepath{app/model/options/core/trees.py} (\codeid{price_asian_mc}, \codeid{price_lookback_mc}, \ldots).
\item MC unifié (European/Bermudan/American via LSMC) : \filepath{app/model/options/service.py} (\codeid{price_option_mc_unified}) appelle \codeid{price_mc_lsmc} dans \filepath{app/model/options/mc_engine.py}.
\item Calibration rough/Volterra : simulateurs dédiés (non-GBM) dans \filepath{app/model/volatility_models/rbergomi/simulator.py} et \filepath{app/model/volatility_models/volterra/simulator.py} utilisent aussi des \codeid{standard_normal}.
\end{itemize}

\section{F. Exemple end-to-end (utilisateur $\to$ calcul)}
\begin{example}[Asian arithmetic (MC) dans l'app]
Dans l'UI Options, un panneau «Asian (arith)» est disponible via \filepath{app/controller/options_controller.py} qui exporte \codeid{price_asian_mc} depuis \filepath{app/model/options/core/trees.py}. Quand l'utilisateur clique «Compute», le flux conceptuel est :
[
(S_0,K,T,\sigma,r,q)\xrightarrow{\text{simulate}} (S^{(i)}_{t_j}) \xrightarrow{\text{moyenne}} \bar S^{(i)} \xrightarrow{\text{payoff}} (\bar S^{(i)}-K)^+ \xrightarrow{\E,\text{disc}} \hat V.
]
\end{example}

\section{G. Limites & extensions crédibles}
\begin{itemize}[leftmargin=*]
\item Le GBM ignore smile/skew : la volatilité implicite observée varie avec $K$ et $T$. D'où la partie calibration/surface.
\item Le MC naïf converge en $O(1/\sqrt{N})$ : variance reduction (antithétiques, control variates) est une extension naturelle (\NonImplem\ explicitement : aucune implémentation dédiée repérée).
\end{itemize}

\begin{jurybox}
Le brownien est le «brique Lego» : il justifie l'apparition de $\sqrt{\Delta t}Z$ dans tous les simulateurs. Dans le code, cette traduction apparaît explicitement (numpy RNG + exponentielle) dans \filepath{app/model/options/core/trees.py}.
\end{jurybox}

% ---------------------------------------------------------
\chapter{Variation quadratique}

\section{A. Problème financier}
Pourquoi une dérivation (Itô) diffère-t-elle du calcul classique ? Parce qu'un brownien a des variations «trop rugueuses» : la somme des carrés des incréments ne tend pas vers 0 mais vers le temps.

\section{B. Théorie}
\begin{definition}[Variation quadratique]
Soit une partition $\pi={0=t_0<\cdots<t_n=T}$. La variation quadratique d'un processus $X$ sur $[0,T]$ est (si la limite existe) :
[
[X]T=\lim{|\pi|\to 0}\sum_{k=0}^{n-1}(X_{t_{k+1}}-X_{t_k})^2.
]
\end{definition}

\begin{proposition}[Cas du brownien]
Pour un brownien $W$, on a $[W]_T=T$ (p.s.).
\end{proposition}

\subsection*{Intuition}
Sur un intervalle $\Delta t$, l'incrément $W_{t+\Delta t}-W_t$ est de taille typique $\sqrt{\Delta t}$. Son carré est donc typiquement $\Delta t$. En sommant $n\simeq T/\Delta t$ termes, on obtient un total $\simeq T$.

\qa
{C'est la raison mathématique de la règle \emph{clé} : $(\dd W)^2=\dd t$ en calcul d'Itô, utilisée pour obtenir les solutions exactes et la PDE de Black--Scholes.}
{Hypothèse : bruit brownien (gaussien, continu).}
{Avec sauts ou volatilité rough, la structure de variation change ; les règles d'Itô se généralisent mais ne sont plus «$(\dd W)^2=\dd t$» au sens classique.}
{Dans l'app, la discrétisation GBM (exponentielle) est précisément choisie pour respecter ces propriétés au premier ordre, cf. \filepath{app/model/options/core/trees.py}.}

\section{C. Dérivations (ultra détaillées)}
\subsection*{Pourquoi $\sum (\Delta W)^2 \to T$}
Pour une partition uniforme $t_k=k\Delta t$, $\Delta W_k=W_{t_{k+1}}-W_{t_k}\sim\Normal(0,\Delta t)$ i.i.d. Donc
[
\E[(\Delta W_k)^2]=\Delta t,\qquad \Var((\Delta W_k)^2)=2(\Delta t)^2.
]
La somme $S_n=\sum_{k=0}^{n-1}(\Delta W_k)^2$ vérifie :
[
\E[S_n]=n\Delta t=T,\qquad \Var(S_n)=n\cdot 2(\Delta t)^2=2T\Delta t\to 0.
]
Donc $S_n\to T$ en $L^2$, donc en probabilité, et (avec un argument plus fin) p.s.

\section{D. Numérique / algorithmes}
En simulation, cela justifie :
\begin{itemize}[leftmargin=*]
\item que l'échelle correcte de l'incrément est \codeid{sqrt(dt)} ;
\item qu'un schéma d'Euler pour $S$ doit inclure les termes d'Itô si on manipule des fonctions non linéaires.
\end{itemize}

\section{E. Implémentation dans le repo}
Même si le code n'écrit pas «$(\dd W)^2=\dd t$», c'est encodé implicitement par :
\begin{itemize}[leftmargin=]
\item \codeid{Z = rng.standard_normal(...)} et \codeid{vol = sigmasqrt(dt)} dans \filepath{app/model/options/core/trees.py}.
\item La construction de processus rough/Volterra où la convolution de noyau est appliquée sur des incréments \codeid{dB = sqrt(dt)*z1} (ex. \filepath{app/model/volatility_models/rbergomi/simulator.py} et \filepath{app/model/volatility_models/volterra/simulator.py}).
\end{itemize}

\section{F. Exemple concret (chiffré)}
\begin{example}[Ordre de grandeur]
Prenons $T=1$ an et $\Delta t=1/252$. Un incrément brownien typique a un écart-type $\sqrt{1/252}\approx 0.063$. Son carré typique vaut $\approx 0.004$. En multipliant par 252 pas, on retombe sur $\approx 1$.
\end{example}

\section{G. Limites & extensions}
\begin{itemize}[leftmargin=*]
\item Si l'on veut une dynamique avec sauts (Merton/Bates) ou rough (rBergomi), la variation quadratique de certains objets change : le code traite ces cas via d'autres modèles (voir Partie II).
\end{itemize}

\begin{jurybox}
La variation quadratique est la justification rigoureuse de la règle de calcul d'Itô. Sans elle, on ne peut pas dériver correctement ni la solution GBM, ni la PDE de Black--Scholes.
\end{jurybox}

% ---------------------------------------------------------
\chapter{Lemme d'Itô (1D et multi-dim)}

\section{A. Problème financier}
Les produits dérivés sont des fonctions $V(t,S_t,\ldots)$ d'états aléatoires. Pour relier :
\begin{itemize}[leftmargin=*]
\item la dynamique du sous-jacent $S_t$,
\item la dynamique du prix de l'option $V_t=V(t,S_t)$,
\end{itemize}
il faut une règle de dérivation adaptée au bruit brownien.

\section{B. Théorie}
\begin{theorem}[Itô 1D]
Soit $X_t$ solution de $\dd X_t=\mu(t,X_t)\dd t+\sigma(t,X_t)\dd W_t$, et $f\in C^{1,2}$. Alors
[
\dd f(t,X_t)=\partial_t f,\dd t+\partial_x f,\dd X_t+\tfrac12,\partial_{xx}f,(\dd X_t)^2,
]
et comme $(\dd X_t)^2=\sigma^2(t,X_t)\dd t$ (via $(\dd W)^2=\dd t$), on obtient
[
\dd f(t,X_t)=\Big(\partial_t f+\mu,\partial_x f+\tfrac12\sigma^2,\partial_{xx}f\Big)\dd t+\sigma,\partial_x f,\dd W_t.
]
\end{theorem}

\begin{theorem}[Itô multi-dim (forme utile)]
Si $\bm{X}_t\in\R^d$ suit $\dd \bm{X}_t=\bm{\mu}\dd t+\Sigma,\dd \bm{W}_t$ et $f\in C^{1,2}$, alors
[
\dd f(t,\bm{X}_t)=\partial_t f,\dd t+\nabla f^\top \dd \bm{X}_t+\tfrac12 \mathrm{Tr}!\big(\Sigma\Sigma^\top \nabla^2 f\big)\dd t.
]
\end{theorem}

\subsection*{Intuition}
Le terme $\tfrac12\sigma^2 f_{xx}\dd t$ est la «correction de convexité» : une fonction convexe appliquée à un bruit a un drift positif (Jensen), et Itô quantifie exactement cet effet à l'ordre $\dd t$.

\qa
{Transformer une dynamique de prix (SDE) en dynamique de prix d'option ; dériver la PDE Black--Scholes ; obtenir des solutions fermées (GBM) ; justifier des schémas numériques.}
{Hypothèses de régularité ($C^{1,2}$) et bruit brownien continu.}
{Si $f$ n'est pas régulière (payoffs avec kink, ex. call $(S-K)^+$), on traite via conditions de viscosité/PDE ou via espérance risque-neutre ; le code utilise souvent des payoffs via \codeid{np.maximum} (non dérivable en $K$) mais ça reste pricer numériquement.}
{Le code encode implicitement Itô via la solution exponentielle de GBM (\filepath{app/model/options/core/trees.py}) et explicitement via la PDE résolue par \codeid{CrankNicolsonBS} (\filepath{app/model/options/engines/pricing.py}).}

\section{C. Dérivations (ultra détaillées)}
\subsection*{Dérivation (Taylor) et apparition de $(\dd W)^2$}
On écrit un développement de Taylor à l'ordre 2 :
[
f(t+\dd t, X_t+\dd X)-f(t,X_t)
= f_t\dd t + f_x\dd X + \tfrac12 f_{xx}(\dd X)^2 + o(\dd t).
]
Or $\dd X=\mu \dd t+\sigma \dd W$. En développant :
[
(\dd X)^2 = (\mu \dd t)^2 + 2\mu\sigma,\dd t,\dd W + \sigma^2 (\dd W)^2.
]
On utilise les ordres de grandeur :
[
\dd t = O(\dd t),\quad \dd W = O(\sqrt{\dd t}),\quad (\dd W)^2=O(\dd t),\quad \dd t,\dd W=O(\dd t^{3/2}),\quad (\dd t)^2=O(\dd t^2).
]
Donc à l'ordre $\dd t$, on néglige $(\dd t)^2$ et $\dd t,\dd W$, mais on \emph{garde} $\sigma^2(\dd W)^2$, et avec $[W]t=t$, on remplace $(\dd W)^2$ par $\dd t$. D'où le terme $\tfrac12\sigma^2 f{xx}\dd t$.

\subsection*{Exemple Itô : $f(x)=\ln x$ et GBM}
Sous GBM : $\dd S=(r-q)S,\dd t+\sigma S,\dd W$.
Pour $f(S)=\ln S$ :
[
f_S=\frac1S,\qquad f_{SS}=-\frac1{S^2}.
]
Itô :
[
\dd \ln S = \frac1S\dd S +\tfrac12\Big(-\frac1{S^2}\Big)(\sigma^2 S^2)\dd t
= (r-q)\dd t+\sigma \dd W - \tfrac12\sigma^2 \dd t
]
soit
[
\dd \ln S = \big(r-q-\tfrac12\sigma^2\big)\dd t + \sigma,\dd W,
]
ce qui justifie la solution exponentielle utilisée en simulation.

\section{D. Numérique / algorithmes}
Deux traductions numériques :
\begin{itemize}[leftmargin=*]
\item \textbf{Solution exacte} (GBM) : exponentielle (stable, $S>0$) $\Rightarrow$ utilisée dans \filepath{app/model/options/core/trees.py}.
\item \textbf{PDE} : si on préfère résoudre la PDE plutôt que simuler, on utilise Crank--Nicolson (\filepath{app/model/options/engines/pricing.py}).
\end{itemize}

\section{E. Implémentation dans le repo}
\begin{itemize}[leftmargin=*]
\item GBM exact : \filepath{app/model/options/core/trees.py}, \codeid{simulate_gbm_paths}.
\item Itô multi-dim (corrélations) : dans rBergomi, construction de $dW=\rho dB + \sqrt{1-\rho^2} dB^\perp$ (voir \filepath{app/model/volatility_models/rbergomi/simulator.py}, variables \codeid{dB}, \codeid{dW}).
\end{itemize}

\section{F. Exemple concret}
\begin{example}[Corrélation en code (rBergomi)]
Dans \filepath{app/model/volatility_models/rbergomi/simulator.py}, on construit deux bruits gaussiens \codeid{z1}, \codeid{z2} puis on corrèle via
[
\dd W = \rho,\dd B + \sqrt{1-\rho^2},\dd B^\perp,
]
ce qui est la traduction directe du fait que $\Corr(\dd W,\dd B)=\rho$.
\end{example}

\section{G. Limites & extensions}
\begin{itemize}[leftmargin=*]
\item Itô suppose continuité : pour les sauts (Merton/Bates) on utilise des formules via fonctions caractéristiques (Partie II) plutôt qu'Itô pur.
\end{itemize}

\begin{jurybox}
Itô est la «règle de dérivation» du monde stochastique. La correction $\tfrac12\sigma^2 f_{xx}$ est \emph{le} point conceptuel à comprendre : c'est elle qui rend GBM log-normal et qui mène à la PDE de pricing.
\end{jurybox}

% ---------------------------------------------------------
\chapter{Risk-neutral pricing, hedging et PDE (Black--Scholes & Feynman--Kac)}

\section{A. Problème financier}
Une option européenne a un payoff $\Phi(S_T)$ à maturité. Objectif : déterminer un prix «sans arbitrage» $V_0$ aujourd'hui, cohérent avec la possibilité de couvrir (hedger) dynamiquement.

\section{B. Théorie}
\subsection*{Définition formelle : absence d'arbitrage et mesure risque-neutre}
Dans un marché frictionless complet (hypothèse BS), absence d'arbitrage $\Rightarrow$ existence d'une mesure $\Qmeas$ telle que le prix actualisé est une martingale :
[
\tilde S_t = e^{-(r-q)t}S_t \text{ est une $\Qmeas$-martingale}.
]
Pour un payoff $\Phi(S_T)$, le prix est :
[
V_0 = e^{-rT}\E^{\Qmeas}[\Phi(S_T)\mid \F_0].
]

\subsection*{Intuition financière}
La mesure $\Qmeas$ n'est pas «la vraie probabilité», c'est une transformation qui fait disparaître les primes de risque dans le drift. Le prix devient une espérance actualisée car on raisonne en \emph{numéraire} (monnaie à taux $r$).

\qa
{Relier couverture (réplication) $\leftrightarrow$ prix ; produire une PDE solvable numériquement ; justifier la formule BS.}
{Marché complet, trading continu, pas de coûts, volatilité constante, taux $r$ et dividend yield $q$ constants.}
{En discret (rebalancement), coûts, sauts, vol stochastique : réplication imparfaite. L'app a des briques «coûteuses/research» (rough/Volterra) qui reconnaissent cette limite.}
{Pricing PDE : \codeid{CrankNicolsonBS} dans \filepath{app/model/options/engines/pricing.py}. Formule BS : \filepath{app/model/options/engines/black_scholes.py} et \filepath{app/model/options/core/greeks.py}.}

\section{C. Dérivations (ultra détaillées)}
\subsection*{Dérivation complète de Black--Scholes via portefeuille couvert}
On suppose sous $\Pmeas$ : $\dd S=\mu S\dd t+\sigma S\dd W$ (le drift réel $\mu$ est inconnu/irrelevant pour le pricing BS).

Soit $V(t,S)$ le prix de l'option. Par Itô :
[
\dd V = V_t \dd t + V_S \dd S + \tfrac12 V_{SS}(\dd S)^2
= \Big(V_t+\mu S V_S+\tfrac12\sigma^2 S^2 V_{SS}\Big)\dd t + \sigma S V_S \dd W.
]
Construisons un portefeuille auto-financé :
[
\Pi = V - \Delta S,
]
où $\Delta$ est choisi pour annuler le terme aléatoire. On calcule :
[
\dd \Pi = \dd V - \Delta,\dd S.
]
En substituant :
[
\dd \Pi = \Big(V_t+\mu S V_S+\tfrac12\sigma^2 S^2 V_{SS}\Big)\dd t + \sigma S V_S \dd W

\Delta(\mu S\dd t+\sigma S\dd W).
]
Choisir $\Delta=V_S$ annule le terme en $\dd W$ :
[
\dd \Pi = \Big(V_t+\tfrac12\sigma^2 S^2 V_{SS}\Big)\dd t.
]
Maintenant, en absence d'arbitrage, un portefeuille sans risque doit rapporter le taux sans risque \emph{sur la valeur investie}. Ici, attention : la position en $S$ doit être financée et il existe un yield $q$ (dividende continu). La forme standard conduit à :
[
\dd \Pi = r\Pi,\dd t = r(V - V_S S)\dd t.
]
En égalant les deux expressions :
[
V_t+\tfrac12\sigma^2 S^2 V_{SS}=r(V - S V_S).
]
En réarrangeant, on obtient la PDE de Black--Scholes avec dividende :
[
\boxed{V_t + (r-q)S V_S + \tfrac12\sigma^2 S^2 V_{SS} - rV = 0.}
]
Condition terminale : $V(T,S)=\Phi(S)$.
\subsection*{Feynman--Kac (PDE $\leftrightarrow$ espérance)}
Le théorème de Feynman--Kac dit essentiellement : la solution de la PDE précédente est l'espérance risque-neutre actualisée. Intuition : on suit la trajectoire stochastique de $S_t$ sous $\Qmeas$ (drift $r-q$) et la martingale du prix actualisé force l'égalité PDE/espérance.

\section{D. Numérique / algorithmes}
\subsection*{Crank--Nicolson : pourquoi ce choix ?}
Crank--Nicolson est un schéma implicite d'ordre 2 en temps (en pratique stable) utilisé pour des PDE paraboliques. Dans l'app, il est encapsulé par la classe \codeid{CrankNicolsonBS}.

\section{E. Implémentation dans le repo}
\subsection*{Solveur PDE}
\begin{itemize}[leftmargin=*]
\item Classe \codeid{CrankNicolsonBS} : \filepath{app/model/options/engines/pricing.py}.
\item La couche service l'utilise directement : \filepath{app/model/options/service.py}, fonction \codeid{compute_price}.
\end{itemize}

\subsection*{Extrait de code (création du solveur depuis le service)}
Dans \filepath{app/model/options/service.py}, \codeid{compute_price} construit le solveur PDE à partir des champs option/marché :

\begin{lstlisting}[caption={Création du solveur PDE : \filepath{app/model/options/service.py}, \codeid{compute_price} (extrait).},label={lst:compute-price}]
cpflag = "p" if cpflag.startswith("p") else "c"
type_raw = str(opt.get("style") or opt.get("exercise") or "Eu").strip().lower()
if type_raw.startswith("am"):
typeflag = "Am"
elif type_raw.startswith("bm"):
typeflag = "Bmd"
else:
typeflag = "Eu"

S0 = float(opt.get("S0") or mkt.get("spot") or 0.0)
K = float(opt.get("strike") or 0.0)
T = float(opt.get("T") or 0.0)
vol = float(opt.get("sigma") or opt.get("iv") or 0.0)
r = float(opt.get("r") or mkt.get("r") or 0.0)
q = float(opt.get("q") or mkt.get("q") or 0.0)

solver = CrankNicolsonBS(typeflag, cpflag, S0=S0, K=K, T=T, vol=vol, r=r, d=q)
price, _, _, _ = solver.CN_option_info()
\end{lstlisting}

\paragraph{Pourquoi ce bout de code est la traduction directe de la PDE ?}
\begin{itemize}[leftmargin=*]
\item \codeid{cpflag} encode la condition terminale call/put $\Phi(S)=(S-K)^+$ ou $(K-S)^+$.
\item \codeid{typeflag} sélectionne les contraintes d'exercice (Eu/Am/Bmd), ce qui correspond à une PDE avec \emph{contrainte d'obstacle} (américaine/bermudane).
\item \codeid{CrankNicolsonBS(...)} construit la grille et résout numériquement la PDE (dans \codeid{CN_option_info}).
\end{itemize}

\section{F. Exemple concret (chiffré)}
\begin{example}[Prix BS (numérique)]
Avec $S_0=100$, $K=100$, $r=2%$, $q=0$, $\sigma=20%$, $T=1$ an, le prix call BS vaut environ $8.916$ (et le put $6.936$).
Ces valeurs sont cohérentes avec la formule implémentée dans \filepath{app/model/options/engines/black_scholes.py} (\codeid{black_scholes_price}).
\end{example}

\section{G. Limites & extensions}
\begin{itemize}[leftmargin=*]
\item Les hypothèses BS sont souvent violées (smile, sauts, vol stochastique) : la Partie II explique comment l'app traite cela via des modèles de surface.
\item En pratique, hedging discret + coûts : la Partie III (RL) propose une approche «décisionnelle», mais le DQN présent est un prototype.
\end{itemize}

\begin{jurybox}
La PDE Black--Scholes vient d'une couverture : choisir $\Delta=V_S$ annule le risque brownien. Le solveur \codeid{CrankNicolsonBS} est l'implémentation numérique de cette PDE, accessible via \codeid{compute_price} dans \filepath{app/model/options/service.py}.
\end{jurybox}

% ---------------------------------------------------------
\chapter{Options vanilles : BS, greeks, implied vol (inversion numérique)}

\section{A. Problème financier}
Dans une salle de marché, on cote souvent une option via sa volatilité implicite $\sigma_{\text{impl}}$ plutôt que son prix. On veut :
\begin{itemize}[leftmargin=*]
\item pricer (BS), obtenir les grecs (sensibilités),
\item inverser BS : prix $\to$ $\sigma_{\text{impl}}$,
\item comparer une surface marché à une surface modèle.
\end{itemize}

\section{B. Théorie}
\subsection*{Formule BS (call/put) avec dividende continu}
Pour $d_1=\frac{\ln(S_0/K)+(r-q+\tfrac12\sigma^2)T}{\sigma\sqrt{T}}$, $d_2=d_1-\sigma\sqrt{T}$ :
[
C = S_0e^{-qT}N(d_1)-Ke^{-rT}N(d_2),\qquad
P = Ke^{-rT}N(-d_2)-S_0e^{-qT}N(-d_1).
]

\subsection*{Grecs (rappel)}
[
\Delta_{\text{call}}=e^{-qT}N(d_1),\quad
\Gamma = \frac{e^{-qT}\varphi(d_1)}{S_0\sigma\sqrt{T}},\quad
\text{Vega}=S_0e^{-qT}\varphi(d_1)\sqrt{T},
]
où $\varphi$ est la densité $\Normal(0,1)$.

\subsection*{Volatilité implicite}
Définition : $\sigma_{\text{impl}}$ est la solution de
[
C_{\text{BS}}(S_0,K,r,q,T,\sigma)=C_{\text{mkt}}.
]
Pour un call européen, $C_{\text{BS}}$ est croissant en $\sigma$ (Vega $>0$), donc l'inversion est un problème 1D bien posé (si le prix respecte les bornes d'arbitrage).

\qa
{Standardiser la comparaison marché/modèle ; construire une surface (K,T) d'IV ; alimenter la calibration (Partie II).}
{Option européenne, modèle BS pour définir l'IV ; prix dans les bornes (pas d'arbitrage statique).}
{IV dépend du modèle : sous sauts/vol stochastique, l'IV résume un modèle plus riche. Pour certains produits (américaines), l'IV BS est un proxy.}
{BS : \filepath{app/model/options/engines/black_scholes.py} et \filepath{app/model/options/core/greeks.py}. Inversion : \filepath{app/model/calibration/implied_vol.py} (\codeid{implied_vol_call}, \codeid{implied_vol_grid}).}

\section{C. Dérivations (ultra détaillées)}
\subsection*{Bornes d'arbitrage et monotonicité}
Pour un call :
[
\max(S_0e^{-qT}-Ke^{-rT},0)\le C \le S_0e^{-qT}.
]
Si $C$ est hors bornes, \emph{aucune} volatilité ne peut l'expliquer dans BS $\Rightarrow$ on renvoie NaN.

\subsection*{Inversion numérique : méthode de Brent}
On résout $f(\sigma)=C_{\text{BS}}(\sigma)-C_{\text{mkt}}=0$.
Brent combine bissection (robuste) et interpolation (rapide) et requiert un encadrement $[\sigma_{\min},\sigma_{\max}]$ tel que $f(\sigma_{\min})f(\sigma_{\max})\le 0$.

\section{D. Numérique / algorithmes}
\begin{itemize}[leftmargin=*]
\item Robustesse : vérifier les bornes d'arbitrage avant de lancer le solveur.
\item Choix d'intervalle : typiquement $\sigma\in[10^{-4},5]$.
\item Complexité : $O(N_{\text{grid}}\times N_{\text{iter}})$ pour une grille d'IV.
\end{itemize}

\section{E. Implémentation dans le repo}
\subsection*{Inversion d'IV (extrait)}
Dans \filepath{app/model/calibration/implied_vol.py} :

\begin{lstlisting}[caption={Inversion IV par Brent : \filepath{app/model/calibration/implied_vol.py}, \codeid{implied_vol_call} (extrait).},label={lst:iv-brent}]
def implied_vol_call(price, S0, K, t, r, q, vol_min=1e-4, vol_max=5.0):
if t <= 0 or S0 <= 0 or K <= 0:
return np.nan
intrinsic = max(S0 * math.exp(-q * t) - K * math.exp(-r * t), 0.0)
upper_bound = S0 * math.exp(-q * t)
if price < intrinsic - 1e-8 or price > upper_bound + 1e-8:
return np.nan

def f(vol):
    return bs_call_price(S0, K, t, r, q, vol) - price

try:
    return float(brentq(f, vol_min, vol_max, maxiter=100, xtol=1e-6))
except Exception:
    return np.nan
\end{lstlisting}

\paragraph{Traduction directe de la théorie.}
\begin{itemize}[leftmargin=*]
\item Les bornes \codeid{intrinsic} et \codeid{upper_bound} implémentent l'absence d'arbitrage statique.
\item \codeid{f(vol)} est $C_{\text{BS}}(\sigma)-C_{\text{mkt}}$.
\item \codeid{brentq} réalise l'inversion numérique robuste.
\end{itemize}

\subsection*{Grecs (proxy vanilla)}
Les grecs sont calculés via BS dans \filepath{app/model/options/core/greeks.py} (\codeid{compute_vanilla_greeks}). Pour des structures multi-jambes, \codeid{compute_option_greeks} prend la première jambe comme proxy (\NonImplem\ au sens «grecs exacts multi-jambes»).

\section{F. Exemple end-to-end}
\begin{example}[Surface modèle Heston $\to$ IV]
Le pricer Heston «V1» (\filepath{app/model/calibration/heston_pricer.py}, \codeid{call_price_cf}) produit des prix, puis \codeid{implied_vol_call} reconstruit l'IV. Par exemple, avec $(\kappa,\theta,\sigma,\rho,v_0)=(2,0.04,0.5,-0.7,0.04)$, $S_0=K=100$, $r=2%$, $q=0$, $T=1$, le code donne un prix $\approx 8.664$ et une IV implicite $\approx 0.194$ (mécanisme : prix $\to$ Brent $\to$ IV).
\end{example}

\section{G. Limites & extensions}
\begin{itemize}[leftmargin=*]
\item L'IV BS est un «résumé» : deux modèles différents peuvent produire la même IV sur un point mais diverger hors grille.
\item Pour les options américaines, l'IV BS n'est pas naturellement définie (il faut un pricer américain en forward). L'app utilise des pricers CRR/LSMC pour ces styles.
\end{itemize}

\begin{jurybox}
La volatilité implicite est l'interface entre marché et modèle. Dans le dépôt, le couple \codeid{bs_call_price}/\codeid{implied_vol_call} (\filepath{app/model/calibration/implied_vol.py}) est la brique d'inversion utilisée pour transformer des prix FFT/MC en surface d'IV calibrée.
\end{jurybox}

% =========================================================
% PARTIE II — Volatility surface & calibration
% =========================================================
\part{Volatility surface & calibration (cœur “quant” de l'app)}

\chapter{Surface d'IV : smile/skew, axes, interpolation, arbitrage}

\section{A. Problème financier}
Le marché cote une \emph{surface} $\sigma_{\text{impl}}(K,T)$ (ou en moneyness $m=K/S_0$). Objectif :
\begin{itemize}[leftmargin=*]
\item représenter proprement les observations (données bruitées et incomplètes),
\item interpoler pour pricer des strikes non cotés,
\item calibrer un modèle paramétrique pour reproduire la surface.
\end{itemize}

\section{B. Théorie}
\subsection*{Définition formelle}
On appelle \emph{surface d'IV} la fonction
[
(m,T)\mapsto \sigma_{\text{impl}}(m,T),
]
où $m=K/S_0$ et $T$ est la maturité en années. Le code de calibration standardise précisément ces deux axes.

\subsection*{Intuition financière (smile/skew)}
\begin{itemize}[leftmargin=*]
\item \textbf{Skew} (actions) : les puts OTM sont chers $\Rightarrow$ IV plus élevée pour $K<S_0$ (demande de protection, levier, crash risk).
\item \textbf{Smile} (FX) : symétrie plus marquée due à la structure de risque différente.
\end{itemize}

\qa
{Construire un objet «marché» qui alimente calibration et pricing ; éviter de pricer sur des trous de données.}
{Suppose qu'on peut représenter le marché par des IV pointwise (ce qui est une convention).}
{Interpolation peut violer non-arbitrage (convexité en strike, monotonicité en maturité). L'app fait une interpolation simple et ne prouve pas l'absence d'arbitrage.}
{La surface est projetée sur une grille fixe via \codeid{build_fixed_grid}/\codeid{make_fixed_grid} dans \filepath{app/model/calibration/market_surface.py}.}

\section{C. Dérivations (ultra détaillées)}
\subsection*{Pourquoi passer en moneyness $m=K/S_0$ ?}
Si le sous-jacent change d'échelle (stock split, variation de prix), une grille en strike absolu devient instable. En moneyness, l'ATM est proche de $m=1$ et on stabilise l'espace de calibration.

\subsection*{Surface discrète et masque d'observation}
On observe des points $(m_i,T_i,\sigma_i)$ irréguliers. Pour calibrer, on construit une grille $(m_j){j=1..M}$, $(T_k){k=1..K}$ et on forme une matrice $\Sigma_{k,j}$.
Comme tous les points ne sont pas observés, on utilise un \emph{masque} $M_{k,j}\in{0,1}$ qui indique quels points proviennent d'observations. Cela permet :
\begin{itemize}[leftmargin=*]
\item de ne pas pénaliser un modèle sur des valeurs interpolées (si on choisit \codeid{fit_to_observed_only=True}),
\item de distinguer «data» vs «remplissage».
\end{itemize}

\section{D. Numérique / algorithmes}
\subsection*{Interpolation dans l'app}
\filepath{app/model/calibration/market_surface.py} :
\begin{itemize}[leftmargin=*]
\item associe chaque point observé au voisin le plus proche sur la grille (remplissage + masque),
\item tente une interpolation \codeid{griddata(method="linear")} si SciPy est disponible,
\item sinon «fallback» par nearest/median.
\end{itemize}

\section{E. Implémentation dans le repo}
\begin{itemize}[leftmargin=*]
\item Grilles par défaut : \codeid{M_GRID_DEFAULT} et \codeid{T_GRID_DEFAULT} dans \filepath{app/model/calibration/market_surface.py}.
\item Construction : \codeid{build_fixed_grid} retourne \codeid{iv_grid}, \codeid{mask}, \codeid{m_grid}, \codeid{t_grid}.
\end{itemize}

\section{F. Exemple end-to-end}
Dans l'onglet «Calibration avancée» (\filepath{app/vue/tabs/tab_advanced_calibration.py}) :
\begin{enumerate}[leftmargin=*]
\item on charge une surface Yahoo (DataFrame avec colonnes \codeid{K,T,S0,iv,type}) via \codeid{CalibrationController.fetch_yahoo_iv_surface} (\filepath{app/controller/calibration_controller.py}) qui appelle \codeid{fetch_iv_surface} (\filepath{app/model/options/data/iv_surface.py}).
\item on canonicalise/filtre dans la Vue, puis on passe au contrôleur une table «surface marché» ;
\item le contrôleur appelle \codeid{make_fixed_grid} pour obtenir (matrice, masque).
\end{enumerate}

\section{G. Limites & extensions}
\begin{itemize}[leftmargin=*]
\item \NonImplem\ : aucune contrainte explicite d'arbitrage statique (convexité en strike) n'est repérée dans les modules de calibration.
\item Extension crédible : calibration sous contraintes (convexité) ou paramétrisations arbitrage-free (SVI pour equity).
\end{itemize}

\begin{jurybox}
La surface n'est pas un «joli plot» : c'est un \textbf{tensor de données} (grille + masque) qui devient l'entrée de la calibration. Dans le code, cette normalisation est concentrée dans \filepath{app/model/calibration/market_surface.py}.
\end{jurybox}